\documentclass[12pt,a4paper]{article}
\usepackage[T1]{fontenc}
\usepackage[utf8]{inputenc}
\usepackage[polish]{babel}
\usepackage{amsmath,amssymb,amsthm,mathtools}
\usepackage[a4paper,margin=2cm]{geometry}
\usepackage{lmodern}
\allowdisplaybreaks

\begin{document}

\begin{center}
{\LARGE Rozwiązania zadań z liczb zespolonych i algebry liniowej}
\end{center}

\section*{Zadanie 1}
\textbf{Działania na liczbach zespolonych w postaci algebraicznej.}

Niech
\[
z_1=a+bi,\qquad z_2=c+di,\qquad a,b,c,d\in\mathbb{R},\ i^2=-1.
\]
Wtedy:
\begin{align*}
z_1+z_2&=(a+c)+(b+d)i,\\
z_1-z_2&=(a-c)+(b-d)i,\\
z_1z_2&=(a+bi)(c+di)=(ac-bd)+(ad+bc)i,\\
\frac{z_1}{z_2}&=\frac{a+bi}{c+di}\cdot\frac{c-di}{c-di}
=\frac{(ac+bd)+(bc-ad)i}{c^2+d^2},\quad (c,d)\neq(0,0).
\end{align*}

\textbf{Pierwiastek kwadratowy liczby zespolonej.}

Szukamy
\[
\sqrt{u+vi}=x+yi,\quad x,y\in\mathbb{R}.
\]
Po podniesieniu do kwadratu:
\[
(x+yi)^2=(x^2-y^2)+2xy\,i=u+vi.
\]
Dostajemy układ:
\[
\begin{cases}
x^2-y^2=u,\\
2xy=v.
\end{cases}
\]
Stąd (po standardowym przekształceniu):
\[
x=\pm\sqrt{\frac{\sqrt{u^2+v^2}+u}{2}},
\qquad
y=\operatorname{sgn}(v)\,\sqrt{\frac{\sqrt{u^2+v^2}-u}{2}},
\]
a drugi pierwiastek to liczba przeciwna.

\bigskip
\textbf{Równanie:}
\[
(2-i)z^2+(7+4i)z+(1+7i)=0.
\]
Dla równania kwadratowego:
\[
a=2-i,\quad b=7+4i,\quad c=1+7i.
\]
\[
\Delta=b^2-4ac.
\]
Liczymy po kolei:
\begin{align*}
b^2&=(7+4i)^2=49+56i+16i^2=33+56i,\\
ac&=(2-i)(1+7i)=2+14i-i-7i^2=9+13i,\\
4ac&=36+52i.
\end{align*}
Zatem
\[
\Delta=(33+56i)-(36+52i)=-3+4i.
\]
Szukamy $\sqrt{\Delta}$:
\[
(1+2i)^2=1+4i+4i^2=-3+4i,
\]
więc
\[
\sqrt{\Delta}=\pm(1+2i).
\]
Wzór na pierwiastki:
\[
z_{1,2}=\frac{-b\pm\sqrt{\Delta}}{2a}=\frac{-(7+4i)\pm(1+2i)}{4-2i}.
\]

\underline{Pierwszy pierwiastek:}
\begin{align*}
z_1&=\frac{-6-2i}{4-2i}\cdot\frac{4+2i}{4+2i}
=\frac{(-6-2i)(4+2i)}{20}
=\frac{-24-12i-8i-4i^2}{20}\\
&=\frac{-20-20i}{20}=-1-i.
\end{align*}

\underline{Drugi pierwiastek:}
\begin{align*}
z_2&=\frac{-8-6i}{4-2i}\cdot\frac{4+2i}{4+2i}
=\frac{(-8-6i)(4+2i)}{20}
=\frac{-32-16i-24i-12i^2}{20}\\
&=\frac{-20-40i}{20}=-1-2i.
\end{align*}

\textbf{Sprawdzenie (dla $z_1=-1-i$):}
\[
z_1^2=(-1-i)^2=2i.
\]
\begin{align*}
(2-i)z_1^2+(7+4i)z_1+(1+7i)
&=(2-i)(2i)+(7+4i)(-1-i)+(1+7i)\\
&=(2+4i)+(-3-11i)+(1+7i)=0.
\end{align*}
Czyli $z_1$ jest poprawnym rozwiązaniem.

\section*{Zadanie 2}
\textbf{Metoda Gaussa.}

Metoda Gaussa polega na:
\begin{itemize}
\item zapisaniu układu jako macierzy rozszerzonej,
\item wykonywaniu dozwolonych operacji elementarnych na wierszach,
\item sprowadzeniu układu do postaci schodkowej (lub schodkowej zredukowanej),
\item odczytaniu rozwiązań (jedno, brak lub nieskończenie wiele).
\end{itemize}

Układ:
\[
\begin{cases}
-7t-3x+5y+2z=-6,\\
5t+x-3y-2z=6,\\
2t-y=2,\\
-2t-2x+2y+z=-1,\\
6t+2x-4y-2z=6.
\end{cases}
\]

Z trzeciego równania:
\[
2t-y=2\quad\Longrightarrow\quad t=1+\frac y2.
\]

Podstawiamy do drugiego:
\begin{align*}
5\left(1+\frac y2\right)+x-3y-2z&=6\\
5+\frac{5y}{2}+x-3y-2z&=6\\
x-\frac y2-2z&=1.
\end{align*}

Podstawiamy do czwartego:
\begin{align*}
-2\left(1+\frac y2\right)-2x+2y+z&=-1\\
-2-y-2x+2y+z&=-1\\
-2x+y+z&=1.
\end{align*}

Podstawiamy do piątego:
\begin{align*}
6\left(1+\frac y2\right)+2x-4y-2z&=6\\
6+3y+2x-4y-2z&=6\\
2x-y-2z&=0\\
x-\frac y2-z&=0.
\end{align*}

Porównując
\[
x-\frac y2-2z=1
\quad\text{oraz}\quad
x-\frac y2-z=0,
\]
od razu dostajemy
\[
z=-1.
\]
Wstawiamy do $x-\frac y2-z=0$:
\[
x-\frac y2+1=0\quad\Longrightarrow\quad x=-1+\frac y2.
\]
A wcześniej mieliśmy
\[
t=1+\frac y2.
\]
Niech parametr wolny
\[
y=\lambda\in\mathbb R.
\]
Wtedy
\[
\boxed{\
(t,x,y,z)=\left(1+\frac\lambda2,\,-1+\frac\lambda2,\,\lambda,\,-1\right),\quad \lambda\in\mathbb R.
}
\]

\textbf{Trzy przykładowe rozwiązania:}
\begin{align*}
\lambda=0&:\quad (t,x,y,z)=(1,-1,0,-1),\\
\lambda=2&:\quad (t,x,y,z)=(2,0,2,-1),\\
\lambda=-2&:\quad (t,x,y,z)=(0,-2,-2,-1).
\end{align*}

\textbf{Sprawdzenie dla }$(1,-1,0,-1)$:
\begin{align*}
-7\cdot1-3(-1)+5\cdot0+2(-1)&=-6,\\
5\cdot1+(-1)-3\cdot0-2(-1)&=6,\\
2\cdot1-0&=2,\\
-2\cdot1-2(-1)+2\cdot0+(-1)&=-1,\\
6\cdot1+2(-1)-4\cdot0-2(-1)&=6.
\end{align*}
Wszystkie równania są spełnione.

\section*{Zadanie 3}
\textbf{Definicja macierzy odwrotnej i warunek istnienia.}

Dla macierzy kwadratowej $A$ macierzą odwrotną nazywamy macierz $A^{-1}$ taką, że
\[
AA^{-1}=A^{-1}A=I.
\]
Macierz odwrotna istnieje wtedy i tylko wtedy, gdy
\[
\det(A)\neq 0.
\]
Wyznaczać ją można np. metodą Gaussa-Jordana z macierzy rozszerzonej $(A\,|\,I)$.

Dla
\[
A=
\begin{bmatrix}
2x-4&x+4&-1\\
2&x+3&-2\\
0&3&-1
\end{bmatrix}
\]
liczymy wyznacznik (rozwinięcie względem 1. wiersza):
\begin{align*}
\det A&=(2x-4)
\begin{vmatrix}
x+3&-2\\3&-1
\end{vmatrix}
-(x+4)
\begin{vmatrix}
2&-2\\0&-1
\end{vmatrix}
+(-1)
\begin{vmatrix}
2&x+3\\0&3
\end{vmatrix}\\
&=(2x-4)\big((x+3)(-1)-(-2)\cdot3\big)
-(x+4)(-2)+(-1)(6)\\
&=(2x-4)(3-x)+2x+8-6\\
&=-2x^2+12x-10\\
&=-2(x^2-6x+5)=-2(x-1)(x-5).
\end{align*}

Zatem $A$ jest odwracalna dla:
\[
\boxed{x\in\mathbb R\setminus\{1,5\}.}
\]

Dla $x=-1$:
\[
A=
\begin{bmatrix}
-6&3&-1\\
2&2&-2\\
0&3&-1
\end{bmatrix},
\qquad
\det A=-24\neq 0,
\]
więc odwrotność istnieje. Otrzymujemy
\[
A^{-1}=
\begin{bmatrix}
-\frac16&0&\frac16\\
-\frac1{12}&-\frac14&\frac7{12}\\
-\frac14&-\frac34&\frac34
\end{bmatrix}.
\]

\textbf{Sprawdzenie:}
\[
AA^{-1}=
\begin{bmatrix}
1&0&0\\
0&1&0\\
0&0&1
\end{bmatrix}=I_3.
\]

\section*{Zadanie 4}
\textbf{Równania prostej i płaszczyzny w $\mathbb R^3$.}

Prosta przez punkt $P_0=(x_0,y_0,z_0)$ o wektorze kierunkowym $\vec v=(a,b,c)$:
\[
\begin{cases}
x=x_0+at,\\
y=y_0+bt,\\
z=z_0+ct,
\end{cases}
\quad t\in\mathbb R,
\]
lub w postaci symetrycznej
\[
\frac{x-x_0}{a}=\frac{y-y_0}{b}=\frac{z-z_0}{c}.
\]

Płaszczyzna o wektorze normalnym $\vec n=(A,B,C)$:
\[
Ax+By+Cz+D=0.
\]

Mamy punkt
\[
P=(4,-1,4)
\]
oraz płaszczyznę
\[
\pi:\ -x+y-z+3=0.
\]
Tu
\[
\vec n=(-1,1,-1),\quad D=3.
\]

Wzór na rzut prostopadły punktu $P$ na płaszczyznę:
\[
H=P-\frac{Ax_0+By_0+Cz_0+D}{A^2+B^2+C^2}\,\vec n.
\]
Liczymy współczynnik:
\begin{align*}
\frac{Ax_0+By_0+Cz_0+D}{A^2+B^2+C^2}
&=\frac{(-1)\cdot4+1\cdot(-1)+(-1)\cdot4+3}{(-1)^2+1^2+(-1)^2}\\
&=\frac{-4-1-4+3}{3}=\frac{-6}{3}=-2.
\end{align*}

Zatem
\[
H=P-(-2)\vec n=P+2\vec n=(4,-1,4)+(-2,2,-2)=(2,1,2).
\]
Punkt symetryczny $P'$ spełnia $H=\frac{P+P'}2$, więc
\[
P'=2H-P=(4,2,4)-(4,-1,4)=(0,3,0).
\]

\[
\boxed{P'=(0,3,0).}
\]

\section*{Zadanie 5}
\textbf{Wzajemne położenie prostych skośnych.}

Proste skośne w $\mathbb R^3$ to proste, które:
\begin{itemize}
\item nie są równoległe,
\item nie przecinają się,
\item nie leżą w jednej płaszczyźnie.
\end{itemize}
Ich odległość to długość najkrótszego odcinka prostopadłego do obu prostych.

Dane:
\[
l_1:\frac{x+3}{-3}=\frac{y-3}{-2}=\frac{z-2}{-2},
\qquad
l_2:\frac{x-2}{4}=\frac{y-3}{2}=\frac{z-3}{4}.
\]

Parametryzacje:
\begin{align*}
l_1&: (x,y,z)=(-3,3,2)+t(-3,-2,-2),\\
l_2&: (x,y,z)=(2,3,3)+s(4,2,4).
\end{align*}

Wektory kierunkowe:
\[
\vec v_1=(-3,-2,-2),\qquad \vec v_2=(4,2,4).
\]
Iloczyn wektorowy:
\[
\vec v_1\times\vec v_2=
\begin{vmatrix}
\mathbf i&\mathbf j&\mathbf k\\
-3&-2&-2\\
4&2&4
\end{vmatrix}
=(-4,4,2),
\quad
\|\vec v_1\times\vec v_2\|=6.
\]

Wektor między punktami bazowymi:
\[
\vec w=(2,3,3)-(-3,3,2)=(5,0,1).
\]
Odległość prostych skośnych:
\[
d=\frac{|\vec w\cdot(\vec v_1\times\vec v_2)|}{\|\vec v_1\times\vec v_2\|}
=\frac{|(5,0,1)\cdot(-4,4,2)|}{6}
=\frac{|-20+0+2|}{6}=\frac{18}{6}=3.
\]

\[
\boxed{d(l_1,l_2)=3.}
\]

\textbf{Punkty realizujące minimalną odległość.}

Niech
\[
A\in l_1:\ A=P_1+t\vec v_1,
\qquad
B\in l_2:\ B=P_2+s\vec v_2,
\]
gdzie
\[
P_1=(-3,3,2),\quad P_2=(2,3,3).
\]
Warunek minimalnej odległości: wektor $\overrightarrow{AB}=B-A$ jest prostopadły do obu kierunków,
\[
(B-A)\cdot\vec v_1=0,
\qquad
(B-A)\cdot\vec v_2=0.
\]
Równoważnie dla $A-B=P_1-P_2+t\vec v_1-s\vec v_2$:
\[
(P_1-P_2+t\vec v_1-s\vec v_2)\cdot\vec v_1=0,
\]
\[
(P_1-P_2+t\vec v_1-s\vec v_2)\cdot\vec v_2=0.
\]
Mamy
\[
P_1-P_2=(-5,0,-1),
\]
\[
(P_1-P_2)\cdot\vec v_1=17,\quad \vec v_1\cdot\vec v_1=17,\quad \vec v_2\cdot\vec v_1=-24,
\]
\[
(P_1-P_2)\cdot\vec v_2=-24,\quad \vec v_2\cdot\vec v_2=36.
\]
Zatem układ:
\[
\begin{cases}
17+17t+24s=0,\\
-24-24t-36s=0.
\end{cases}
\]
Po uproszczeniu:
\[
\begin{cases}
17t+24s=-17,\\
2t+3s=-2.
\end{cases}
\]
Rozwiązanie:
\[
s=0,\qquad t=-1.
\]
Stąd
\[
A=P_1-\vec v_1=(-3,3,2)+(3,2,2)=(0,5,4),
\]
\[
B=P_2=(2,3,3).
\]
Długość:
\[
|AB|=\|(2,-2,-1)\|=\sqrt{4+4+1}=3,
\]
zgodnie z obliczoną odległością.

\[
\boxed{A=(0,5,4),\quad B=(2,3,3).}
\]

\section*{Zadanie 6}
\textbf{Diagonalizacja macierzy.}

Macierz $A$ jest diagonalizowalna, gdy istnieje macierz odwracalna $P$ i diagonalna $D$ takie, że
\[
A=PDP^{-1}.
\]
Kolumny $P$ to wektory własne macierzy $A$, a na przekątnej $D$ stoją odpowiadające im wartości własne.

Dla
\[
A=
\begin{bmatrix}
2&2\\
3&3
\end{bmatrix}
\]
wyznaczamy wartości własne:
\[
\det(A-\lambda I)=
\begin{vmatrix}
2-\lambda&2\\
3&3-\lambda
\end{vmatrix}
=(2-\lambda)(3-\lambda)-6
=\lambda^2-5\lambda
=\lambda(\lambda-5).
\]
Zatem
\[
\lambda_1=5,\qquad \lambda_2=0.
\]

Dla $\lambda_1=5$:
\[
(A-5I)v=0
\Rightarrow
\begin{bmatrix}
-3&2\\
3&-2
\end{bmatrix}
\begin{bmatrix}x\\y\end{bmatrix}=0
\Rightarrow -3x+2y=0.
\]
Możemy wziąć wektor własny
\[
v_1=\begin{bmatrix}2\\3\end{bmatrix}.
\]

Dla $\lambda_2=0$:
\[
Av=0
\Rightarrow
\begin{bmatrix}
2&2\\
3&3
\end{bmatrix}
\begin{bmatrix}x\\y\end{bmatrix}=0
\Rightarrow x+y=0.
\]
Możemy wziąć
\[
v_2=\begin{bmatrix}1\\-1\end{bmatrix}.
\]

Tworzymy
\[
P=\begin{bmatrix}2&1\\3&-1\end{bmatrix},
\qquad
D=\begin{bmatrix}5&0\\0&0\end{bmatrix}.
\]
\[
\det P=2\cdot(-1)-1\cdot3=-5\neq 0,
\]
więc $P$ jest odwracalna i
\[
P^{-1}=\frac1{-5}
\begin{bmatrix}
-1&-1\\
-3&2
\end{bmatrix}
=
\begin{bmatrix}
\frac15&\frac15\\
\frac35&-\frac25
\end{bmatrix}.
\]

\textbf{Sprawdzenie:}
\[
PD=
\begin{bmatrix}
2&1\\3&-1
\end{bmatrix}
\begin{bmatrix}
5&0\\0&0
\end{bmatrix}
=
\begin{bmatrix}
10&0\\15&0
\end{bmatrix}.
\]
Dalej:
\[
PDP^{-1}=
\begin{bmatrix}
10&0\\15&0
\end{bmatrix}
\begin{bmatrix}
\frac15&\frac15\\
\frac35&-\frac25
\end{bmatrix}
=
\begin{bmatrix}
2&2\\3&3
\end{bmatrix}=A.
\]
Diagonalizacja wykonana poprawnie.

\end{document}
